% !TeX root = ../regression_logistique.tex
\chapter{Rappels de prérequis}
\label{ann:prerequis}

\orient{%
réactiver les outils de statistique, d'analyse et de géométrie employés par le
cours}{%
aucun}{%
référence ponctuelle ; chaque section indique le chapitre ou l'annexe de retour}

\section{Résumer une colonne de nombres}

Soit $v_1,\dots,v_m$ les valeurs observées dans une colonne.

\textbf{Moyenne.} $\mu=\frac{1}{m}\sum_i v_i$. Elle se déplace avec chaque
valeur : ajouter mille points à la plus haute note déplace $\mu$ de
$1000/m$.

\textbf{Médiane.} La valeur qui coupe la colonne triée en deux moitiés de même
effectif ; pour un effectif pair, la demi-somme des deux valeurs centrales.
Déplacer la plus haute note ne la change pas tant que l'ordre des valeurs
centrales est préservé. C'est cette insensibilité qu'on appelle
\emph{robustesse}.

\textbf{Écart-type.} $\sigma=\sqrt{\frac{1}{m}\sum_i(v_i-\mu)^2}$ : la racine de
la moyenne des carrés des écarts à la moyenne. Il s'exprime dans l'unité de la
colonne et mesure la dispersion typique autour de $\mu$.

\textbf{Cote standard.} $(v-\mu)/\sigma$ exprime un écart en nombre d'écarts-types.
La colonne transformée a pour moyenne $0$ et pour écart-type $1$.

Employés au chapitre~\ref{ch:standardisation}.

\begin{rappel}[Valeur manquante]
Une case vide représente une valeur inconnue. Trois traitements existent :
écarter la ligne,
écarter la colonne, ou attribuer une valeur --- l'\emph{imputation}. Le choix
dépend du mécanisme qui a produit l'absence, question développée
dans~\cite{littlerubin2019} et que ce cours n'aborde pas.
\end{rappel}

\section{Exponentielle, logarithme, cote}

$e^{a}$ est strictement positive, strictement croissante, et transforme les
sommes en produits : $e^{a+b}=e^{a}e^{b}$. Elle parcourt $]0,+\infty[$ quand $a$
parcourt $\R$.

$\log$ est sa réciproque : définie sur $]0,+\infty[$, strictement croissante,
elle parcourt $\R$. Trois égalités servent constamment :
\[
  \log(ab)=\log a+\log b,\qquad
  \log\frac{a}{b}=\log a-\log b,\qquad
  \log a^{b}=b\log a .
\]
Avec $\log 1=0$ et $\log e^{a}=a$.

Le logarithme employé dans tout le document est le logarithme népérien, celui de
base $e$ ; c'est celui que \texttt{math.log} calcule en Python.

\textbf{Cote.} Le rapport $p/(1-p)$ d'une probabilité à son complément. Une
probabilité de $0{,}8$ donne une cote de $4$ : quatre fois plus de chances que
l'événement se produise que le contraire. Quand $p$ parcourt $]0,1[$, la cote
parcourt $]0,+\infty[$ et son logarithme parcourt $\R$.

Employés aux chapitres~\ref{ch:nuage} et~\ref{ch:erreur}, et à
l'\annexeref{ann:logit}.

\section{Vecteurs, produit scalaire, norme}

Dans $\R^n$ muni d'un repère orthonormé, le \emph{produit scalaire} de
$u=(u_1,\dots,u_n)$ et $v=(v_1,\dots,v_n)$ vaut
\[
  u\cdot v=\sum_{j=1}^{n} u_jv_j = \|u\|\,\|v\|\cos\theta,
\]
où $\theta$ est l'angle entre les deux vecteurs et
$\|u\|=\sqrt{\sum_j u_j^2}$ la \emph{norme} de $u$, c'est-à-dire sa longueur.

Trois conséquences employées au chapitre~\ref{ch:nuage} :

\begin{panelitemize}[accent]
  \panelbullet{$u\cdot v=0$ équivaut à $u\perp v$, les deux vecteurs étant non
    nuls ;}
  \panelbullet{$\|v\|\cos\theta$ est la longueur de la projection de $v$ sur la
    direction de $u$, comptée avec un signe ;}
  \panelbullet{multiplier $u$ par $\lambda>0$ multiplie $u\cdot v$ par
    $\lambda$ sans changer le signe ni la direction.}
\end{panelitemize}

\begin{vigilance}[Le repère doit être orthonormé]
Angle, perpendiculaire et distance n'ont le sens usuel que si les deux axes
portent la même unité. Sur les notes brutes, où un axe se gradue en points et
l'autre en centaines de points, ces notions dépendent du barème. La
standardisation du chapitre~\ref{ch:standardisation} fournit un système d'unités
commun ; elle ne rend pas pour autant la géométrie obtenue équivalente à celle
des données brutes, point discuté au même chapitre.
\end{vigilance}

\section{Dérivée, dérivée partielle, gradient}

La \emph{dérivée} $f'(a)$ mesure la variation de $f$ par unité de variation de
son argument, au voisinage de $a$ : c'est la pente de la tangente en ce point.
Son signe indique le sens de variation, sa valeur absolue la rapidité.

Trois règles suffisent aux calculs du document :
\[
  (u+v)'=u'+v',\qquad
  \left(\frac{1}{u}\right)'=-\frac{u'}{u^{2}},\qquad
  \bigl(f\circ g\bigr)'(a)=f'\bigl(g(a)\bigr)\,g'(a) .
\]
La troisième est la \emph{règle de chaîne}, employée dans toute
l'\annexeref{ann:gradient}.

La \emph{dérivée partielle} $\partial J/\partial w_j$ est la dérivée de $J$
considérée comme fonction du seul $w_j$, les autres coefficients étant tenus
pour constants. Le \emph{gradient} $\nabla J$ est le vecteur de ces dérivées
partielles.

Deux propriétés servent au chapitre~\ref{ch:descente} :

\begin{panelitemize}[rulecolor]
  \panelbullet{$-\nabla J$ est la direction de plus forte décroissance locale de
    $J$ au point considéré ;}
  \panelbullet{$\nabla J$ est orthogonal à la ligne de niveau de $J$ passant par
    ce point.}
\end{panelitemize}

\begin{rappel}[Approximation locale]
Pour un déplacement $\delta$ petit,
$J(\w+\delta)\approx J(\w)+\nabla J(\w)\cdot\delta$. Cette approximation est
locale : elle cesse de décrire $J$ dès que $\delta$ sort du voisinage où la
pente mesurée reste représentative. C'est la raison pour laquelle un pas de
descente trop long peut faire monter le coût
(ch.~\ref{ch:descente}).
\end{rappel}

\section{Dérivée seconde et courbure}

La dérivée seconde $f''$ mesure la variation de la pente. Positive sur un
intervalle, elle y rend $f'$ croissante, donc $f'$ y traverse zéro au plus une
fois.

Pour une fonction de plusieurs variables, la \emph{hessienne} est le tableau
carré des dérivées secondes $\partial^2 J/\partial w_j\partial w_k$. La courbure
de $J$ dans une direction $\mathbf v$ vaut $\mathbf v^{\!\top}H\mathbf v$. Ces
notions servent à l'\annexeref{ann:convexite}.
